\subsection{MaterialSwap event}
\label{sec:verification:material-swap-event}

This test isolates the \texttt{MaterialSwap} MPM event.  A stationary elastic block is generated in \texttt{ParticleRegion1}.  A matching, initially empty destination subregion is generated for \texttt{ParticleRegion2}.  At $t_s=0.5$ the event transfers the active particles from the source region to the destination region.

\begin{figure}[h]
  \centering
  \includegraphics[width=0.70\linewidth]{\CaseSourceDir/setup_material_swap.png}
  \caption{MaterialSwap setup.  The source block initially occupies \texttt{ParticleRegion1}; a zero-volume destination object creates the matched destination subregion without initial particles.  The center tracer is used to check the material-type step and stationary kinematics.}
\end{figure}

There is no applied load, so the ideal kinematic solution is
\begin{equation}
  \mathbf{x}(t)=\mathbf{x}_0, \qquad \mathbf{v}(t)=\mathbf{0}.
\end{equation}
The event-specific ideal solution is the tracer material-type step
\begin{equation}
  m(t)=\begin{cases}
  0, & t < t_s,\\
  1, & t \ge t_s.
  \end{cases}
\end{equation}
The primary metric is the material-type mismatch count away from the event-output coincidence window.  Secondary metrics are maximum tracer displacement and maximum velocity norm.

\IfFileExists{\CaseOutputDir/material_swap_results.tex}{\input{\CaseOutputDir/material_swap_results.tex}}{\paragraph{Generated results.} Results are produced by \texttt{postProcess\_materialSwap.py} after the GEOS job finishes.}
